\documentclass[11pt,a4paper]{article}
\usepackage[T1]{fontenc}
\usepackage{lmodern}
\usepackage[margin=1in]{geometry}
\usepackage{amsmath,amssymb,amsthm}
\usepackage{hyperref}
\usepackage{enumitem}
\hypersetup{colorlinks=true,urlcolor=blue,linkcolor=blue}
\newcommand{\Gr}{\operatorname{Gr}}
\newcommand{\Res}{\operatorname{Res}}
\newcommand{\dd}{\,d}
\title{Positive Geometry: Canonical Forms, Positivity, and the Amplituhedron\\[4pt]\large Expository Learning Notes}
\author{Kumar Azad}
\date{9 November 2026}
\begin{document}
\maketitle
\begin{abstract}
These expository notes develop the elementary mathematical foundation
of positive geometry through boundary residues and canonical forms.
Worked examples include the oriented interval and standard triangle,
followed by the positive Grassmannian and the tree amplituhedron map.
The material summarizes established literature and makes no claim of
original research.
\end{abstract}
\tableofcontents

\section{Definition and boundary recursion}
A positive geometry is, under suitable technical hypotheses,
a pair $(X,X_{\geq0})$, where $X$ is a complex projective
algebraic variety defined over $\mathbb R$ and $X_{\geq0}$
is a suitably oriented closed semialgebraic region of its real points.
It admits a unique nonzero rational top-degree form $\Omega(X,X_{\geq0})$
having only logarithmic singularities on boundaries, such that
\begin{equation}
\boxed{\Res_{B}\Omega(X,X_{\geq0})=\Omega(B,B_{\geq0}).}
\label{eq:recursion}
\end{equation}
Locally near a boundary specified by $t=0$, the polar part has the form
\begin{equation}
\Omega=\frac{dt}{t}\wedge\Omega_B+
\text{terms regular along }t=0.
\end{equation}
The induced orientation controls signs. An oriented point has
canonical form $+1$ or $-1$. See \cite{ABL} for the precise definition.

\section{An oriented interval}
Let $a<b$, with the interval oriented by increasing $x$.
We seek a one-form with only simple poles at $a$ and $b$.
Try
\begin{equation}
\Omega_{[a,b]}=\frac{A\,dx}{(x-a)(b-x)}.
\end{equation}
The complex residue at $x=a$ equals $A/(b-a)$.
Normalizing it to $+1$ gives $A=b-a$. Therefore
\begin{equation}
\begin{aligned}
\boxed{\Omega_{[a,b]}(x)}
&=\frac{(b-a)\,dx}{(x-a)(b-x)}\\
&=\frac{dx}{x-a}+\frac{dx}{b-x}\\
&=d\log\left(\frac{x-a}{b-x}\right).
\end{aligned}
\label{eq:interval}
\end{equation}
The residues at $x=a$ and $x=b$ are $+1$ and $-1$,
respectively, in the usual complex-residue convention.
There is no additional pole at infinity.

\section{The standard triangle}
Define the standard simplex by
\[
\Delta=\{(x,y)\in\mathbb R^2:x\geq0,\quad
y\geq0,\quad x+y\leq1\}.
\]
Choose ambient orientation $dx\wedge dy$.
A rational two-form with logarithmic poles on all three facets
is
\begin{equation}
\boxed{\Omega_\Delta=
\frac{dx\wedge dy}{x\,y\,(1-x-y)}.}
\label{eq:triangle}
\end{equation}
The boundary $x=0$ provides the normalization:
\[
\Omega_\Delta=\frac{dx}{x}\wedge
\frac{dy}{y(1-x-y)},\qquad
\Res_{x=0}\Omega_\Delta=\frac{dy}{y(1-y)}.
\]
This is exactly the interval form on that edge.
Along the boundary $y=0$, rearranging the wedge gives
\[
\Omega_\Delta=\frac{dy}{y}\wedge
\left(-\frac{dx}{x(1-x-y)}\right),\qquad
\Res_{y=0}\Omega_\Delta=-\frac{dx}{x(1-x)}.
\]
The sign is consistent with the induced boundary orientation.
Canonical forms on polygons can also be evaluated by
triangulation: the poles on interior diagonals cancel in
an oriented sum \cite{ABL}.

\section{The positive Grassmannian}
The real Grassmannian $\Gr(k,n)$ parameterizes $k$-planes
in $\mathbb R^n$.
A point is represented by the row span of a full-rank
$k\times n$ matrix $C$, up to invertible row transformations.
For ordered columns $I=\{i_1<\cdots<i_k\}$, define the Pl\"ucker
coordinate (maximal minor)
\begin{equation}
\Delta_I(C)=\det(C_{i_1},\ldots,C_{i_k}).
\end{equation}
The \emph{positive} Grassmannian $\Gr_{>0}(k,n)$ consists of
points admitting a representative with every ordered maximal
minor strictly positive. Its closure is the \emph{nonnegative}
Grassmannian $\Gr_{\geq0}(k,n)$. The real dimension is
$\dim\Gr(k,n)=k(n-k)$.

\subsection{A $2\times3$ example}
Consider
\[
C=\begin{pmatrix}1&1&1\\0&1&2\end{pmatrix}.
\]
The ordered maximal minors are
\[
(\Delta_{12},\Delta_{13},\Delta_{23})=(1,2,1)>0,
\]
so $[C]\in\Gr_{>0}(2,3)$.

\subsection{A Pl\"ucker relation}
For $[C]\in\Gr(2,4)$, the coordinates obey the identity
\begin{equation}
\Delta_{12}\Delta_{34}
-\Delta_{13}\Delta_{24}
+\Delta_{14}\Delta_{23}=0.
\label{eq:plucker}
\end{equation}
This polynomial relation expresses the fact that the maximal
minors are constrained, rather than independent coordinates.
The boundary-cell structure of $\Gr_{\geq0}(k,n)$ is described
by positroid cells \cite{Postnikov}.

\section{The tree amplituhedron}
Fix $n\geq k+m$ and an $n\times(k+m)$ matrix $Z$
whose ordered maximal minors are strictly positive.
For $[C]\in\Gr_{\geq0}(k,n)$, define
\begin{equation}
Y=\operatorname{rowspan}(CZ)\in\Gr(k,k+m).
\end{equation}
The tree amplituhedron is the image
\begin{equation}
\boxed{\mathcal A_{n,k,m}(Z)=
\left\{\operatorname{rowspan}(CZ):
[C]\in\Gr_{\geq0}(k,n)\right\}.}
\label{eq:amplituhedron}
\end{equation}
When $k=1$ and $m=2$, positive external data lead
to a convex polygon in projective two-space, illustrating the idea
in elementary terms. The $m=4$ construction is relevant
to planar $\mathcal N=4$ supersymmetric Yang--Mills theory.
The relationship between canonical forms and scattering
amplitudes is more subtle than identifying a form with a
physical amplitude without additional kinematic and
supersymmetric steps \cite{AT,ABCGPT}.
This special framework is not a general formula for QCD amplitudes.

\section{Exercises}
\begin{enumerate}[label=\arabic*.]
\item Re-derive \eqref{eq:interval} from the endpoint residues.
\item Compute all three triangle facet residues and verify
the orientation signs.
\item Calculate the maximal minors of a general $2\times4$
matrix and verify \eqref{eq:plucker}.
\item Compare the dimension $k(n-k)$ of $\Gr(k,n)$ with
the dimension $km$ of the target $\Gr(k,k+m)$.
\end{enumerate}

\begin{thebibliography}{9}
\bibitem{ABL} N. Arkani-Hamed, Y. Bai, T. Lam,
\emph{Positive Geometries and Canonical Forms},
JHEP 11 (2017) 039.
\url{https://arxiv.org/abs/1703.04541}
\bibitem{Postnikov} A. Postnikov,
\emph{Total positivity, Grassmannians, and networks},
2006. \url{https://arxiv.org/abs/math/0609764}
\bibitem{AT} N. Arkani-Hamed, J. Trnka,
\emph{The Amplituhedron},
JHEP 10 (2014) 030.
\url{https://arxiv.org/abs/1312.2007}
\bibitem{ABCGPT} N. Arkani-Hamed, J.~L. Bourjaily,
F. Cachazo, A.~B. Goncharov, A. Postnikov, J. Trnka,
\emph{Scattering Amplitudes and the Positive Grassmannian}.
\url{https://arxiv.org/abs/1212.5605}
\end{thebibliography}
\end{document}
